\documentclass[12pt]{article}
\usepackage{amsmath,amssymb,geometry,setspace,url}
\geometry{margin=1in}
\setstretch{1.2}
\begin{document}
\title{SKANN-SSL --- Full-Scale System: Diagnostics, Deployment \& Deliverables}
\author{Oravont Systems LLP}
\date{}
\maketitle

\section*{Document F (Very Short Version)}
Full-Scale SKANN-SSL System: Diagnostics, Export, Deployment, Deliverables

\section*{1. Purpose}
Document F provides a concise summary of the diagnostics, deployment workflow, dataset strategy, and deliverables for the full-scale SKANN-SSL system. It complements Documents A–E.

The objectives are:

Define evaluation criteria for learned acoustic embeddings.

Specify model export and deployment requirements.

Describe dataset integration for large-scale self-supervised learning.

List the final deliverables for the SKANN project.

Explain how Mini-SKANN transitions into the full SKANN system.

\section*{2. Diagnostics Suite}
\section*{2.1 Embedding Quality Diagnostics}
Redundancy Index: Measures the magnitude of off-diagonal elements of the cross-correlation matrix  used in Barlow Twins self-supervised learning. Lower values indicate better decorrelation.

Dimensional Utilisation Score: Fraction of embedding dimensions whose variance exceeds a threshold, indicating whether the embedding space is well-used.

Temporal Stability: Cosine similarity between embeddings of overlapping waveform segments to ensure consistent encoding.

\section*{2.2 Robustness Diagnostics}
Augmentation Stress Test: Evaluates the invariance of embeddings under time shifts, gain jitter, masking, and band dropout, the same perturbations used in Barlow Twins SSL training.

SNR Robustness Sweep: Tests embedding performance across SNR values from  dB to  dB.

Propagation and Doppler Variation: Applies simple propagation filters or Doppler stretch/compression to verify encoder robustness to underwater channel effects.

\section*{2.3 Clustering Diagnostics}
Cluster Persistence: Checks the stability of HDBSCAN cluster labels across runs. HDBSCAN is a density-based clustering method that does not require a pre-selected cluster count.

Cluster Separability: Uses silhouette scores and distance ratios to evaluate cluster quality.

Cluster Signature Extraction: Computes the average embedding, waveform, and spectrum for each cluster to support interpretation.

3. Export and Deployment Plan

3.1 ONNX Export

The encoder must export cleanly to ONNX (Open Neural Network Exchange), a portable format that enables inference on CPUs, GPUs, and embedded devices.

Modules to export:

SKConv1D front-end (Selective Kernel 1D Convolution: a multi-kernel convolution block that adaptively fuses features using attention-like weights)

1D SKA backbone

SKConv2D backbone (Selective Kernel 2D Convolution, used after reshaping the features into a time–frequency-like representation)

Hybrid encoder fusion layer

Projection head removed for inference

Input: , Output embedding:

\section*{3.2 Deployment Targets}
ARM Cortex-A processors

DSP architectures

Edge GPUs (e.g., NVIDIA Jetson Nano or Orin)

Shore-side compute nodes

\section*{3.3 Inference Optimisation}
Optional 8-bit quantisation

Fusing convolution–normalisation–activation

Use of TensorRT (NVIDIA’s high-performance inference optimiser) for accelerated deployment on Jetson-class devices

\section*{4. Dataset Integration}
4.1 Two-Phase SSL Strategy

Phase 1: Synthetic Pretraining Uses physics-inspired vessel signatures, broadband propulsion noise, cavitation bursts, and Doppler modifications to initialise the encoder.

Phase 2: Real-World SSL

Datasets include:

NOAA: US National Oceanic and Atmospheric Administration’s passive acoustic archives.

MBARI: Monterey Bay Aquarium Research Institute hydrophone data.

JAMSTEC: Japan Agency for Marine-Earth Science and Technology observatory recordings.

DCLDE: Detection, Classification, Localisation, and Density Estimation Workshop datasets.

\section*{4.2 Standardisation Requirements}
Sampling rate  Hz

Clip length  samples

RMS normalisation

Metadata fields: source, timestamp, depth, SNR tag

\section*{4.3 Dataset Expansion}
Multichannel HAVS acoustic vector-sensor arrays

Long-term seabed monitoring stations

Fleet-noise signature archives for unsupervised vessel characterisation

\section*{5. Deliverables}
5.1 Code and Model Deliverables

Implementation of Stages  to

SSL-trained encoder weights

ONNX deployment model

Normalisation parameters

Augmentation configuration files

\section*{5.2 Documentation Deliverables}
\section*{Document A: Underwater Acoustics}
\section*{Document B: System Overview}
Document C: DSP and Sampling Conventions

\section*{Document D: Selective Kernel Network Architecture}
Document E: Encoder and SSL Pipeline

Document F: Diagnostics and Deployment (this document)

5.3 Visual and Analytical Outputs

Embedding visualisation methods such as:

PCA: Principal Component Analysis (linear dimensionality reduction)

UMAP: Non-linear manifold learning for cluster visualisation

TSNE: Non-linear technique for revealing cluster structure

Additional outputs:

Cluster heatmaps and silhouette diagrams

Robustness plots (SNR, augmentation, Doppler)

Block diagrams for deployment architecture

ONNX export package

6. Scaling Path: Mini-SKANN to Full SKANN

Mini-SKANN provides:

Basic synthetic generator

Minimal SKConv1D / SKConv2D blocks

Simple Barlow Twins SSL (a redundancy-reduction SSL method)

Basic clustering

Full SKANN adds:

Advanced physics-based synthetic generator

Deeper SKA networks with dilation

Multi-layer SKConv2D hierarchy

Hybrid SSL using Barlow Twins and VICReg (Variance–Invariance–Covariance Regularisation)

Multi-sensor HAVS fusion

Deployment-grade optimisation with TensorRT/ONNX Runtime

\section*{7. Conclusion}
Document F describes the diagnostics, deployment, dataset requirements, and scaling strategy needed to operationalise the full SKANN-SSL system.

Appendix A: List of Symbols

: Sampling frequency (Hz)

: Number of samples per 1-second clip

: Embedding dimension

: Cross-correlation matrix used in SSL

SNR: Signal-to-noise ratio

\section*{Appendix B: Short Forms}
NOAA: National Oceanic and Atmospheric Administration

MBARI: Monterey Bay Aquarium Research Institute

JAMSTEC: Japan Agency for Marine-Earth Science and Technology

DCLDE: Detection, Classification, Localisation, Density Estimation Workshop

ONNX: Open Neural Network Exchange

PCA: Principal Component Analysis

UMAP: Uniform Manifold Approximation and Projection

TSNE: t-distributed Stochastic Neighbour Embedding

VICReg: Variance-Invariance-Covariance Regularisation

SKConv1D / SKConv2D: Selective Kernel Convolution Layers

TensorRT: NVIDIA inference optimisation runtime

\end{document}